
\subsection{Theory Directions}
In our research were able to highlight several directions that one could take.
We mainly highlight to directions: \textbf{HF-PPAD} in  \textbf{PPAD}
and topological representations of directed graphs

\subsubsection{\textbf{HF-PPAD} in  \textbf{PPAD}}

In our sections \ref{sec:pc-plus-eol} and \ref{sec:sc-excess}, made references to
how Kleene logic can be incroportated with \textbf{PPAD} problems. For
the \textit{nD-StrongSperner} problems we observed that \textit{Kleene logic}
provides a natural extension and is equivalent to its non hazard counterpart.
For other problems such as \textit{StrongSperner}, we observed that this
is not clear by current complexity assumptions and indeed proving \textbf{PPAD}-hardness
would be essential for establishing the counting hardness of \textit{PureCircuit}.
But one could ask further questions such as, could we define
complexity classes such as \textbf{HF-PPAD} and could we show that
$\textbf{HF-PPAD} = \textbf{PPAD}$.
Along similar line, we can also ask the questions as to existence of problems
\textit{HF-nD-Counting-SS} and if they are parsimonious to their non hazard-free variants.
We know that the non counting version can be resolved parsimoniously
but we argue that it is not so simple when we omit panchromatic cubes as by doubling
the cubes, we may introduce more solutions.


Even more naturally, one can also ask further questions such as, how does hazard-freeness
affect \textbf{TFNP} as a whole. We conjecture that one can make
similar observations to the \textit{StrongTucker} problem \cite{aisenberg_2DTuckerPPA_2020},
in order to create similar classes such as \textbf{HF-PPA} as
it also uses a very similar colouring argument as the \textit{StrongSperner} argument.
Morevore we were able to observe how we can use \textit{Kleene logic} with
\textit{PLS} subproblems like \textit{Tarski}. We believe that this shift
in perspective may also give insights to other areas such as the circuit complexity \textit{Hazard-Free Circuits}.


\subsubsection{Parsimonious Topological Representations of \textbf{PPAD} problems }


From the previous sections we observed how to reduce \textit{SourceOrExcess}$(2,1)$ to a topological problem.
Just like in the \textit{EndOfLine} problem, we found it essential to reduce it to a \textit{StrongSperner} problem
with polynomial sides and very high dimensionality. We believe that one could abuse the \textit{Snake Embedding} methodology
that we proposed earlier but there are two main questions that need to be resolved:
\begin{enumerate}
    \item Will the number of solutions between foldings stay the same?
    \item Is there a natural way to define \textit{Counting-SS}?
\end{enumerate}
More interestingly, we made references of $\alpha-\textit{SUnbalanced}(\Delta_i, \Delta_o)$
and how they can affect the counting nature of each other. Like we mentioned previously, for
\textit{2D-StrongSperner} and \textit{1-MC}, we cannot guarantee parsimonious reductions, but
when we mapped our instances to a $3D$ space that was indeed possible. We therefore make the following conjecture:

\begin{conjecturebox}{\textit{nD-CountingSS} and \textit{Unbalanced}}{}
    There exist some polynomial computable function $f :\mathbb{N}^3 \to \mathbb{N}$, such that
    given $\alpha,\Delta_i, \Delta_o \in \mathbb{N}$ and $m= f(\alpha, \Delta_i, \Delta_o)$
    $$
        \textbf{PPAD}(\alpha\textit{-SUnbalanced}(\Delta_i, \Delta_o)) \subseteq
        \textbf{PPAD}(m\textit{D-CountingSS})
    $$
\end{conjecturebox}
What the above conjecture essentially describes is, given a number of known sources and the
maximum number of input/output vertices of graph, can we create some \textit{nD-CountingSS}
that is parsimonious to its graphical representation? Since we know that the problem \textit{Unbalanced}
can parsimoniously count every variant, can we also reduce \textit{Unbalanced} to a topological
problem? A reduction like that would be an important milestone for the usage topological \textit{PureCircuit} constructions
for enumerating solutions.


\subsubsection{PureCircuit Gate Extension} %% TODO 
The \textit{PureCircuit} instance we defined uses the gate set
\textit{NOR, Purify}. In fact, Deligkas et al. \cite{deligkas_PureCircuitTightInapproximability_2024}
demonstrated that if $G$ is our gate set, then for the problem to be \textbf{PPAD}-complete,
it must be the case that
\begin{gather*}
G \supseteq \{\textsc{Purify}, v\} \text{ for } v \in \{\textsc{Nor}, \textsc{Nand}\} \\
\text{ or }\\
G \supseteq \{\textsc{Purify}, \textsc{Not}, v\}  \text{ for } \{\textsc{And}, \textsc{Or}\} \\
\end{gather*}
The proof for the \textbf{PPAD}-inclusion can be summarised as expressing a \textit{PureCircuit} instances
as a \textit{Brouwer} function from $[0,1]^n \to [0,1]^n$, given fixed point $x^\star$ then
we denote the folloiwng assignment: if $x^\star_i \in \{0,1\}$ then $v_i = x^{\star}_i$, otherwise
$v_i = \bot$. We argue that there may exist gates which exhibit any behaviour of the sort $f \in [0,1]^k \to [0,1]^m$
such that $f$ is contiuous, where they may make our reductions easier and more straight forward.
This could raise additional directions as to the behaviour of the additional functions
and whether they may affect the combinatorial nature of the \textit{PureCircuit} problem.

\subsection{Software future Directions}

\subsubsection{Algorithmic Extensions}
Although the aforementioned algorithms work as expected, there are several improvements that can be made.
For all three of the algoritmic methods we mentioned, the first and most important extension
is to treat weakly directly connected components individually. Trivially an assignment in one
weakly connected component should not affect the solution set of the other one.
We can therefore reduce the search spaces of our problems by finding
the solution sets in each graph seperately. Moreover for the
\textit{backtracking} algorithm, we can extend the above idea by considering
the cartesian product of the solution sets, meaning:
let $G_i$ denote a (weakly) connected component of $G$ and $\textbf{sol}(\cdot)$
denotes its solution set. Its trivial to observe that
$$
    \textbf{sol}(G) = \bigtimes_{k \in [k]}  \textbf{sol}(G_k)
$$

We can also optimise the backtracking algorithm even further.
First one can observe that if our graph was acyclic, then
the permutation of the nodes with no \textit{in-degree} solutions essentially
dictate the assignments of the downstream nodes (excluding nodes incident
to purify gates). With this in mind, we can apply the following sequence
of steps to make the backtracking algorithm simpler:
\begin{enumerate}
    \item Identify the strongly connected components (SCCs) of the circuit, using polynomial
          time algorithms such as Kosaraju's Algorithm \cite{sharir_StrongconnectivityAlgorithmIts_1981} 
          or Gabow's Algorithm \cite{gabow_PathbasedDepthfirstSearch_2000}. 
    \item Treating SCCs as an individual node, identify nodes by their distances from
          base nodes (nodes with no incoming edges)
    \item Extend the backtracking algorithm so instead of the using the is start label,
          use the distance from base instead. The greater distance a node has from the
          base, the less impact it has and therefore should be postponed in later stages of the
          backtracking calls.
\end{enumerate}

We argue that the above heuristic may result in improved average times but of course
further research can be done in order to determine more suitable selection heuristics.
Furthermore, an additional heuristic that we failed to consider is using the notion
of alternation or varying path.
This has been introduced by \cite{eichelberger_HazardDetectionCombinational_1965}
and the idea is the ``application'' of the gates to propagate values.
Their algorithm can be summarised as such:
Start with an initial set of nodes, and apply gates of each circuit enough times.
For nodes that alternate their values between $0,1$, replace their value with $\bot$
Repeat until no altrenation occurs. This algorithm is still exponential
in its running time but one could deduce useful insights as to notion
of alternating paths.
It has to be noted that it may not be entirely clear how such construction
could be applied as the purify gate has a non-deterministic output, meaning
for the $\bot$ value, the outputs could be multiple.
